\chapter{体积、质量、密度}\label{ch:g7:volume-mass-density}

去年你用克除以立方厘米，识破了一位假冒的金匠。
今年，数学课给了你一项新本领：用字母表示数。
密度正是施展它的好地方——一个简短的公式，
就装下了三种算法、一套侦探式的工具，
以及整条漂浮规律，
全都写进了五个符号里。

\section{公式}

\begin{definition}[公式形式的密度]\label{def:g7:volume-mass-density:formula}
物质的\emph{密度}\index{密度} $\rho$（希腊字母
\emph{rho}）是指物质单位体积的质量：对于
质量为 $m$、体积为 $V$ 的样品，
\[
\rho = \frac{m}{V}.
\]
若 $m$ 以克为单位，$V$ 以立方厘米为单位，则 $\rho$ 的单位为
\emph{克每立方厘米}（\unit{g/cm^3}）——正是去年的
“每立方厘米的质量”，如今换上了统一的名称。
水的密度为 $\rho = \qty{1.0}{g/cm^3}$。由于一升等于
$1000$ 立方厘米，同一个数也可以读作每升千克数：
水为每升 \qty{1.0}{kg}。
\end{definition}

\begin{example}[读懂字母]\label{ex:g7:volume-mass-density:reading}
公式是一句简写的话。$\rho = m/V$ 的意思是：“求密度，
就把样品的质量除以它的体积。”字母只是
占位符——对于\emph{这一块}鹅卵石，$m$ 取
\qty{75}{g}，$V$ 取 \qty{30}{cm^3}：
\[
\rho = \frac{m}{V} = \frac{75}{30} = \qty{2.5}{g/cm^3}.
\]
与去年相同的除法，但这个公式记住了
以后每一块鹅卵石的算法。
\end{example}

\begin{proposition}[一个公式的三种算法]\label{prop:g7:volume-mass-density:recipes}
由于密度把质量与体积按比例联系起来，一个公式
可同时满足三种需求：
\begin{enumerate}
  \item \emph{求密度}：$\rho = m/V$——用质量除以
        体积；
  \item \emph{求质量}：$m = \rho \times V$——每立方厘米
        含 $\rho$ 克，一共 $V$ 个
        立方厘米；
  \item \emph{求体积}：$V = m/\rho$——把质量按每份 $\rho$ 克分
        份，份数就是
        立方厘米数。
\end{enumerate}
\end{proposition}

\begin{example}[算法的应用]\label{ex:g7:volume-mass-density:atwork}
\emph{由体积求质量}：\qty{250}{cm^3} 的铁
（$\rho = 7.9$）质量是多少？$m = 7.9 \times 250 = \qty{1975}{g}$——
差不多两千克，大约相当于一只咖啡杯的体积。
\emph{由质量求体积}：\qty{540}{g} 的橡木
（$\rho = 0.6$）占多大体积？$V = 540 \div 0.6 = \qty{900}{cm^3}$。
\emph{合理性检验}：铁的答案应当重而小，
木头的答案应当轻而大——密度表是
判断步骤的最好帮手。
\end{example}

\begin{omfigure}
\begin{tikzpicture}
\begin{axis}[omaxis, width=10.8cm, height=6.2cm, xbar,
    xlabel={密度（\unit{g/cm^3}）},
    symbolic y coords={cork, oak, ice, water, aluminium, iron,
      copper, gold},
    ytick=data, yticklabels={软木, 橡木, 冰, 水, 铝, 铁, 铜, 金}, xmin=0, xmax=21,
    bar width=7pt, xtick={0,1,5,10,15,20},
    nodes near coords, nodes near coords style={font=\footnotesize},
    every axis plot/.append style={fill=omDef, fill opacity=0.55,
      draw=omDef}]
  \addplot coordinates {(0.2,cork) (0.6,oak) (0.9,ice) (1.0,water)
    (2.7,aluminium) (7.9,iron) (9.0,copper) (19.3,gold)};
\end{axis}
\end{tikzpicture}

{\small 密度表宛如一条天际线：从软木轻如羽毛的
$0.2$，到黄金雄伟的 $19.3$。水的 $1.0$ 是漂浮的
分界线。}
\end{omfigure}

\section{测量密度}

\begin{method}[液体的密度]\label{met:g7:volume-mass-density:liquid}
液体不能单独放在天平盘上称量——
因此借用去皮的办法：
\begin{enumerate}
  \item 把空量筒放在天平上，并将其去皮
        调零；
  \item 倒入一个便于计算的体积——例如 $V = \qty{80}{cm^3}$ 的
        待测液体，并正确读取弯月面；
  \item 此时天平显示的正是液体本身的质量，例如
        $m = \qty{68}{g}$；
  \item 相除：$\rho = 68 \div 80 = \qty{0.85}{g/cm^3}$——
        比水轻，很可能是某种油。
\end{enumerate}
\end{method}

\begin{example}[固体的密度：从头到尾]\label{ex:g7:volume-mass-density:solid}
一枚奖牌声称是银制的（$\rho = 10.5$）。称量：
$m = \qty{84}{g}$。量筒读数：水面从 \qty{60.0}{cm^3} 升至
\qty{68.0}{cm^3}，因此 $V = \qty{8.0}{cm^3}$。于是
$\rho = 84 \div 8.0 = \qty{10.5}{g/cm^3}$——密度吻合，
它很可能确实是银制的。（若是用银镀层包裹的廉价金属，
这里就会露馅，除非精心挑选，
正如周末问题将要展示的那样。）
\end{example}

\begin{remark}[公式中的单位]\label{rem:g7:volume-mass-density:units}
公式只有单位一致才是可靠的。若 $m$ 以克计，
$V$ 以 \unit{cm^3} 计，则 $\rho$ 以 \unit{g/cm^3} 计；
若改用千克和升，结果为 \unit{kg/L}——幸好，
对任何物质，数值都相同。但若把克与升混用，
公式会毫无怨言地算出荒谬的结果。专业人士的准则是：
计算之前，先列出并核对单位；
计算之后，把单位写进答案。
\end{remark}

\section{密度的思考}

\begin{example}[合金与中间情形]\label{ex:g7:volume-mass-density:alloys}
将两种物质混合，混合物的密度落在\emph{两者之间}，
且更接近含量较多的那一种。青铜（铜 $9.0$ 加少量锡）
的密度约为 $8.8$；王冠的金银合金介于 $10.5$ 与 $19.3$ 之间；
海水（即溶有盐的水）的密度约升至 $1.03$。
“介于之间”本身就是线索：
一块标称“纯金”的物体若读数为 $12$，
便是掺入了别种物质的供认。
\end{example}

\begin{example}[密度决定漂浮的世界]\label{ex:g7:volume-mass-density:floating}
漂浮规则现在换上正式装束：物体在液体中漂浮，
当 $\rho_{\text{物体}} < \rho_{\text{液体}}$ 时。橡木（$0.6$）
放在水（$1.0$）中：漂浮。冰（$0.9$）放在油（$0.85$）中：下沉——
请查表，不要凭直觉。游泳者（$\approx 1.0$）
置身于著名的超高盐度湖（$\approx 1.2$）中：像软木塞一样漂浮，
手里还拿着一份报纸。一个不等式，管住了整座港湾。
\end{example}

\begin{remark}[注意整体，而非物质本身]\label{rem:g7:volume-mass-density:parcel}
$\rho = m/V$ 描述的是\emph{物质}；船舶和充气的救生衣
则是\emph{物体整体}——物质加上封闭的空气——
而参与漂浮规则比较的，是整体的 $m/V$。
公式两者皆可处理，只要代入正确的 $m$ 和
$V$：对物质，用钢材本身的密度；对整体，用整个船体的
质量除以整个船体的体积。大多数关于漂浮的“悖论”，
其实只是两种计算方式被混为一谈。
\end{remark}

\section{练习}

\begin{exercise}[$\star$]\label{exo:g7:volume-mass-density:1}
写出密度公式，说明每个字母的含义，并给出
与克和立方厘米相配的单位。
\end{exercise}

\begin{exercise}[$\star$]\label{exo:g7:volume-mass-density:2}
计算密度：$m = \qty{270}{g}$，$V = \qty{100}{cm^3}$。
这是天际线图中的哪种金属？
\end{exercise}

\begin{exercise}[$\star$]\label{exo:g7:volume-mass-density:3}
请用正确的公式：求 \qty{40}{cm^3} 铜的质量
（$\rho = 9.0$）；\ 求 \qty{1930}{g} 金的体积
（$\rho = 19.3$）。
\end{exercise}

\begin{exercise}[$\star$]\label{exo:g7:volume-mass-density:4}
在\cref{met:g7:volume-mass-density:liquid}中，为何要先
去皮？随后哪两个读数代入公式？
\end{exercise}

\begin{exercise}[$\star$]\label{exo:g7:volume-mass-density:5}
某液体：$V = \qty{50}{cm^3}$，$m = \qty{70}{g}$。其密度是多少？
冰块（密度为 $0.9$）放入其中会漂浮还是下沉？
\end{exercise}

\begin{exercise}[$\star$]\label{exo:g7:volume-mass-density:6}
为什么水的密度以 \unit{g/cm^3} 表示，
或以千克每升表示时数值相同？
见\cref{rem:g7:volume-mass-density:units}，该处对克与升混用
提出了怎样的警示？
\end{exercise}

\begin{exercise}[$\star$]\label{exo:g7:volume-mass-density:7}
一只标有“纯金”的手镯，$m = \qty{58}{g}$，
$V = \qty{5.0}{cm^3}$。求其密度，并
以密度表为依据给出判断。
\end{exercise}

\begin{exercise}[$\star\star$]\label{exo:g7:volume-mass-density:8}
一立方米是一个 $100 \times 100 \times 100$
厘米的立方体。它等于多少 \unit{cm^3}——以及一立方米水的质量
是多少千克？（答案
可解释为何水床会让房东发愁。）
\end{exercise}

\begin{exercise}[$\star\star$]\label{exo:g7:volume-mass-density:9}
一座“银”质奖杯：$m = \qty{420}{g}$，放入量筒后水面从
\qty{200}{cm^3} 升至 \qty{260}{cm^3}。求其密度，
它落在表中哪两种物质之间？
“介于之间”悄悄说明了什么？
\end{exercise}

\begin{exercise}[$\star\star$]\label{exo:g7:volume-mass-density:10}
空瓶质量为 \qty{380}{g}；注满至其 \qty{75}{cL}
刻度线的某种未知液体后，总质量为 \qty{1010}{g}。求
该液体以 \unit{g/cm^3} 表示的密度，并推断其种类。
\end{exercise}

\begin{exercise}[$\star\star$]\label{exo:g7:volume-mass-density:11}
一只空心铝制浮标，总体积为 \qty{4000}{cm^3}，
总质量为 \qty{1200}{g}。求该\emph{整体}的
密度，并预测它会漂浮还是下沉——然后借助
\cref{rem:g7:volume-mass-density:parcel}，说明为何铝本身的
$2.7$ 不是应当查用的正确数值。
\end{exercise}

\begin{exercise}[$\star\star\star$]\label{exo:g7:volume-mass-density:12}
设计一套完整的方案，判断一条链条是否为纯铜：
列出每项测量、所用的配方、纯度的预期数值，以及\emph{两}条诚实的理由，
说明你的结论仍可能出错（考虑空心的链环，
以及密度恰好接近铜的狡猾混合物）。
还需要什么额外的检验，
才能使结论更加可靠？
\end{exercise}

\section{问题：废品场里的侦探}

\begin{problem}[{周末问题 --- 废品场金属侦探的一个下午；
四批待鉴定的金属，一块冒充的锭；
密度测定的局限}]\label{pb:g7:volume-mass-density:1}
废品场按金属的\emph{本质}而非外表收购金属，
场里的侦探手持天平、一只大号量筒，
还有一张密度表。你是学徒。
（表格摘录：铝 $2.7$；铁 $7.9$；铜 $9.0$；铅 $11.3$；
金 $19.3$；钨 $19.3$。）

\medskip\noindent\textbf{第一部分 --- 四批金属。}
\begin{enumerate}
  \item A 批，一块灰色锭：$m = \qty{5400}{g}$，排开水的体积
        $V = \qty{2000}{cm^3}$。求密度并判断种类。
  \item B 批，一卷金属丝：$m = \qty{1800}{g}$，
        $V = \qty{200}{cm^3}$。求密度并判断种类。
  \item C 批，一块暗淡而沉重的板：$m = \qty{4520}{g}$，
        $V = \qty{400}{cm^3}$。求密度并判断种类——并说明
        为何 C 批必须戴手套、慎重对待？
  \item D 批，一袋混杂的浅色废料：$m = \qty{8100}{g}$，
        $V = \qty{1500}{cm^3}$。证明 D 批的密度落在
        表中两种金属之间，并说明这袋废料
        最可能含有什么。
\end{enumerate}

\medskip\noindent\textbf{第二部分 --- 价格与预测。}
废品场按质量付款，但按体积规划运输。
\begin{enumerate}[resume]
  \item 一位买家想从 A 批中切下 \qty{540}{g} 铝。
        这需要多大体积的锭？
  \item 将 B 批的铜线熔化成立方块，每块体积为
        \qty{25}{cm^3}。求每块的质量，并求
        这卷铜线能整块地做出多少块？
  \item 一只货箱最多可载 \qty{20}{kg}。它最多可合法地装
        多少块这样的铜立方块？
  \item 卡车的车槽还可再装 \qty{3000}{cm^3}
        的铁废料。侦探被允许
        多装多少质量（以千克计）？
\end{enumerate}

\medskip\noindent\textbf{第三部分 --- 冒充者。}
一位卖家带来一块闪亮的“金”锭：
$m = \qty{3860}{g}$，排开水的体积 $V = \qty{200}{cm^3}$。
\begin{enumerate}[resume]
  \item 求该锭的密度。它是否符合金的
        特征？
  \item 侦探不为所动，查看表中的最后一行，
        叹了口气。哪一种更便宜的金属
        恰好具有与金相同的密度？
        这对密度测定的局限说明了什么？
  \item 废品场的两件仪器中，哪一件在此被骗过了：
        天平、量筒，还是两者，抑或都没有？请准确说明
        密度测定\emph{确实}诚实地证实了该锭的
        什么。
  \item 从本课程前几章中，提出一项进一步的物理检验，
        能在不损坏该锭的前提下区分金与它的仿制品。
        （它们的熔点相差悬殊——金约为 \qty{1064}{\celsius}，
        仿制品则远高于此处的每一座熔炉——
        但废品场不会熔化一块可能是宝物的东西：
        请找出更温和的证据，例如磁铁对铁芯的判定，
        或敲击金属棒的声音，
        并诚实地为你的选择辩护。）
\end{enumerate}

\medskip\noindent\textbf{第四部分 --- 侦探的技艺。}
\begin{enumerate}[resume]
  \item 写出侦探的三句信条：密度能证明什么，
        只能暗示什么，
        以及单凭它永远不能做什么。
\end{enumerate}
\end{problem}
